\section{One finite probability space for all layer lengths}\label{sec:finite}

Lemma~\ref{lem:final-gates} assigns random bits to keys and to sets of keys
in a directed acyclic graph. A term of Section~\ref{sec:numerical} can have
arbitrarily many distinct ancestors. Assigning a fresh bit to each ancestor
would make the probability space depend on the term. We instead index the
bits by finite descriptions of keys, called names. Two keys with the same
name receive the same bits. We choose the depth of the descriptions at
random. At each fixed key, with high probability the equalities between
names in its neighborhood are exactly the equalities of keys in some legal
graph. This local guarantee will suffice both for the occurrence probability
and for the path in Lemma~\ref{lem:primary}.

\subsection{Recognizing a legal neighborhood}

Fix the parameters of the gate construction. The \emph{unfolding} at a key
to radius \(R\) is the rooted tree of its move sequences of length at most
\(R\). A named unfolding records the node data, the available labeled move
slots and an equivalence relation on the occurrences, namely equality of
names. At each fixed depth there are only finitely many such patterns up
to renaming, since branching, slots and node data are bounded.

A color refinement is a function of a finite pattern from earlier rounds.
Give the function an arbitrary default on patterns where it was undefined,
proceeding through the rounds in order. Extend the local prescriptions of
raw, barrier and screen tests to all patterns in the same manner. These
extensions use only equality of names and preserve the bit indices on
realizable patterns.

\begin{lemma}[Local realizability]\label{lem:local-extension}
For every \(R\) there exists \(R'\ge R\) such that the radius-\(R\)
restriction of a named unfolding to radius \(R'\) is realized in the graph
model of Section~\ref{sec:probability} if it passes the following checks.
\begin{enumerate}
\item Equal names have the same node data. Whenever a move slot is
exposed at two occurrences of one name, it has the same target name.
\item The basic move rules and the prefix and witness rules for flags hold
wherever the neighborhoods needed to test them are exposed. In these
checks the colors are calculated by the extended refinement rules.
\item Identifying equal names in the graph of exposed moves creates no
directed cycle.
\end{enumerate}
\end{lemma}
\begin{proof}
If no such radius exists, there are counterexamples of arbitrarily large
radii. By successive extraction, stabilize their patterns at each finite
depth. In the limiting tree identify equivalent occurrences. The first
check makes each move well defined on the quotient. A directed cycle
would use finitely many exposed moves and name equalities, contradicting
the third check at a finite depth.

The second check gives the basic moves. Suppose the first rounds of
colors and flags already satisfy the rules of the graph model. These
rounds admit a legal completion by making all subsequent flags false.
The next refinement reads only the completed earlier rounds, so its
extended value is its specified value on the graph. The second check now
gives the prefix and witness conditions for the next flags. Induction
produces a legal graph, possibly infinite, whose root realizes the
stabilized radius-\(R\) pattern. This contradicts the choice of the
counterexamples.
\end{proof}

Put \(d=D_{h-1}\). Choose \(R\) large enough to evaluate the color
refinements and gate tests at the root, including the flags and gates at
charged keys along every path in the final screen. Such a path uses at
most \(dL_0\) moves. Fix the resulting \(R'\) from
Lemma~\ref{lem:local-extension}. A pattern passing its checks inherits
the probability bounds and the missed path of Lemma~\ref{lem:final-gates}.
Indeed, realization preserves every equality among the indices read by
these tests, so the local bits have precisely their graph-model joint law.

\subsection{Giving names to the keys of a term}

The two blank layers of each site are its shores. A nonleaf blank term has
its two parent steps as basic moves. Attach two temporary blank leaves to
each micro leaf, one on each shore of its site. From a micro term, a
straight run of parent steps in either direction now reaches a first blank
term, its direction key. Compress this straight run, or \emph{leg}, into
one basic move to its key. A witness jump in the same direction will end
at a strictly older micro term on the leg, with the same direction key.
When we return to the original terms we stop a path at the micro leaf
preceding a temporary blank.

Choose a radius \(j\) uniformly among the integers
\(N_*,\ldots,2N_*-1\). In round zero this range is a single bin. In each
later round subdivide every preceding bin into equally many consecutive
bins of equal size. We first choose the subdivision counts, then take
\(N_*\) sufficiently large and divisible by the number of final bins.
Each bin in a round has its own flags and witness slots.

To define names, retain the moves and flags for \emph{all} bins, the site,
the node type and an auxiliary symbol. The \(j\)-name of a term is the
tree of its moves truncated after \(j\) levels, with these data through
depth \(j\). It records neither identities of terms nor the positions of
micro layers within a site. Once the parameters and subdivision counts
are fixed, the number of outgoing move slots is bounded. Choose a finite
ordered alphabet larger than the number of strict ancestors that can be
reached in \(2N_*\) enriched moves. Each term takes the first symbol
absent at those ancestors. The symbol is part of its name, but the local
rules use names only for equality. In particular, the check radius
\(R'\) was fixed before choosing the auxiliary alphabet.

Process terms in order of construction, with temporary blanks first and
parents before their merges.
\begin{itemize}
\item A blank term has known outgoing moves. Assign its symbol, form
its names and compute every round's refinement outputs for every \(j\).
Use equality of \(j\)-names and the bins containing \(j\) in the
previous rounds. Call a round's bin \emph{stable at the key} if the vector
of outputs for all input colors is constant on that bin. Round zero is
stable.
\item In a bin a micro term is \emph{live} if both direction keys are
stable in that bin and in every containing bin of earlier rounds. Its
calculated color is then constant on the bin. On each leg set flags for
exactly the colors of the live micro terms encountered, including the
source. Witness each flag by the first such occurrence of its color.
If the source is the witness, no jump is needed. Otherwise the
corresponding slot jumps to the witness. Assign the source its auxiliary
symbol after all these moves have been determined.
\end{itemize}
In the second step the source's liveness and color are calculated from
its two keys, both strictly older, so the specification is recursive.
Every witness keeps the same direction key along the leg. Its liveness
in the preceding bin supplies the flag for its prefix color.

\begin{lemma}[Finite names]\label{lem:finite-names}
Given \(\varepsilon>0\), the subdivisions and \(N_*\) can be chosen
independently of the numbers of micro layers so that, for each fixed blank
term, with probability at least \(1-\varepsilon\) over \(j\),
\begin{enumerate}
\item the bin containing \(j\) is stable at this key in every round, and
\item the unfolding at this key computed with the bins of \(j\) and
equality of \(j\)-names passes the radius-\(R'\) checks of
Lemma~\ref{lem:local-extension}.
\end{enumerate}
\end{lemma}
\begin{proof}
Inside a bin of the preceding round, the flags and move slots read by the
current refinement are fixed. It inspects at most \(u\) occurrences. Here
\(u\) may be chosen uniformly for the fixed local rules, since the
refinement reads one bin per earlier round. Increasing \(j\) only refines
the partition of these occurrences by name equality. Each strict
refinement adds at least one class. Between refinements the output vector
is constant, so at most \(u-1\) of the new bins can be unstable. Choose
enough bins at each subdivision that, summed over the finitely many
rounds, the probability of instability is less than \(\varepsilon/2\).

With these counts fixed, the radius-\(R'\) tree using the slots of
\emph{all} bins contains at most \(M\) occurrences, independently of
\(N_*\) and the lengths of the sites. Its partition by \(j\)-names
can strictly refine at most \(M-1\) times between depths \(j\) and
\(j+1\). For every other \(j\), equal \(j\)-names among these occurrences
have identical \((j+1)\)-names and hence, in every exposed slot, targets
with identical \(j\)-names. This gives the first check. The decoration
gives the second, since each witness is evaluated with the same bin and
the same name equalities that supplied its flag.

For the third check, suppose the identified graph has a simple directed
cycle of length \(m\le M\). Follow the cycle's slots in the actual tree
starting at its first representative. After one step the subtree agrees
with the next representative to depth at least \(j-1\): the chosen slot
occurs in the starting name and its target has the next representative's
name. At each succeeding step at most one more level of agreement is
lost. If \(N_*>M\), the terminal term agrees with the starting one at
least through depth zero. It therefore has the same auxiliary symbol,
although it is a strict ancestor within \(m\le 2N_*\) enriched moves.
The choice of symbols excludes this. Finally choose \(N_*\) so large
that \((M-1)/N_*<\varepsilon/2\). This bounds the probability of failure
of the checks and completes the proof.
\end{proof}

\begin{proof}[Proof of the Primary Lemma (Lemma~\ref{lem:primary})]
Call a blank key \emph{good} at \(j\) when both conclusions of
Lemma~\ref{lem:finite-names} hold there, with
\(\varepsilon=L_0^{-10}\). An outcome selects \(j\), a uniform final color,
and a common table of independent gate bits, independently of one another.
Supply entries for every syntactically possible index of a raw or barrier
bit. Replace keys by names and sets of keys by sets of names, and retain
the prescribed stages, colors, and other index data, including the sizes
of sets when these affect a bit's law. The screens read this same table.
All trees naming the bits have bounded depth over a finite alphabet.
Hence the outcome space \(\Omega\) is finite and works for every choice
of the numbers of micro layers.

For a micro term \(v\) let \(P(v)\) consist of the outcomes for which
both direction keys are good, the term has the selected final color,
and both keys' final gates succeed in that color. For a blank term put
\(P(v)=\varnothing\).

The two keys of \(v\) are good together with probability at least
\(1-2L_0^{-10}\). Fix such a radius. Conditional on selecting the color
of \(v\), their joint raw successes have probability at least \(q^{-4}\)
by Lemma~\ref{lem:final-gates}, since their opposite shores distinguish
their names. Condition on these successes. At each key the sum of the
conditional failure probabilities of its screens is \(o(1)\). To see
why the conditioning at the other key is allowed, if its name occurs
in the local test, it denotes a key of the same site in the realizing
graph. If its name is absent, its raw bits are independent of that test.

At each key the required barrier bits pass with probability at least
\(\exp(-4h\lambda)\). Shared barrier bits require the same success, so
the probability that both prescriptions pass is at least
\(\exp(-8h\lambda)\). These bits are independent of the raw bits.
Subtracting the sum of the conditional screen failure probabilities
therefore gives, with \(k_f=a^{o(1)}\) final colors,
\[
 \Pr(P(v))\ \ge\ (1-2L_0^{-10})\,k_f^{-1}q^{-4}
       \bigl(\exp(-8h\lambda)-o(1)\bigr)\ \ge\ q^{-5}.
\]
The last inequality holds for sufficiently large \(a\), since \(q=a^w\)
and \(h,\lambda,w\) depend only on \(C\). The allowed powers of two \(a\)
in~\eqref{eq:prob-parameters} give arbitrarily large powers of two \(q\).

Fix an outcome and a starting blank term in site \(s\). Say a charge
succeeds if its key is good and its final gate succeeds in the selected
color. Expand each compressed move to its straight sequence of parent
steps. A potential occurrence on this sequence, with the source included
and the destination excluded, forces a successful charge at that move.
Indeed, such an occurrence is a live micro term of the selected color
on the source's leg, and the common direction key is good with
a successful gate. The color at that occurrence sets the source's flag.
The half-open convention includes a micro leaf whose move ends at a
temporary blank.

Follow the basic spine pointing out of \(s\). If no charge succeeds, it
ends at a leaf without an occurrence. Otherwise its first successful
charge takes the basic move to its key in a site \(s'\ne s\). Its
expansion so far can have occurrences only on this last leg, which lies
in \(s'\). At this good key local realizability and Lemma~\ref{lem:final-gates}
supply a continuation to a leaf on which all charges to successful final
gates of the selected color belong to \(s'\). The path lifts by following
its labeled slots in the actual term. Our choice of \(R\) includes every
move of the path and every gate test at its charged keys. Thus every
potential occurrence on its expanded path belongs to that site.
The two pieces use at most \(3L_0\) and \(dL_0\) compressed moves
respectively. Each expands in a single direction, so the path has at
most \((d+3)L_0\) turns. Stop at the original leaf before any temporary
blank. This gives the path and the uniform turn bound of
Lemma~\ref{lem:primary}.
\end{proof}
